% !TEX root = ../../main.tex
\subsection{功能测试}

本节对 FoodShield 系统的核心功能进行正向测试，覆盖用户注册、PID 生成、口令认证、订单创建与 Token 生成、Token 验证、骑手接单、双向实时通信、消息加密存储、Merkle 快照生成、条件溯源和三端一致性验证等关键功能。测试结果如表~\ref{tab:func_test}所示，14 项测试全部通过。系统各端运行界面截图详见第四章 4.10 节。

\begin{table}[H]
\centering
\caption{功能测试用例与结果}
\label{tab:func_test}
\footnotesize
\setlength{\tabcolsep}{4pt}
\begin{tabular}{clllcc}
\hline
\textbf{编号} & \textbf{测试项} & \textbf{输入描述} & \textbf{预期结果} & \textbf{实际} & \textbf{结论} \\
\hline
F01 & 用户注册 & 用户名+密码+验证码 & 生成PID & 通过 & $\checkmark$ \\
F02 & PID生成 & 注册后自动 & 前端不展示 & 通过 & $\checkmark$ \\
F03 & 用户登录 & 用户名+密码 & SM3验证通过 & 通过 & $\checkmark$ \\
F04 & 创建订单 & 登录用户 & 生成orderID+Token & 通过 & $\checkmark$ \\
F05 & Token验证 & 合法Token & valid=true & 通过 & $\checkmark$ \\
F06 & Token过期 & 超时时间戳 & valid=false & 通过 & $\checkmark$ \\
F07 & 骑手注册 & 骑手信息 & 生成rider\_pid & 通过 & $\checkmark$ \\
F08 & 骑手接单列表 & 骑手登录 & 不展示PID & 通过 & $\checkmark$ \\
F09 & 骑手接单 & 骑手+orderID & 状态变更 & 通过 & $\checkmark$ \\
F10 & 实时通信 & 同订单房间 & 消息收发 & 通过 & $\checkmark$ \\
F11 & 加密存储 & 发送消息 & content为密文 & 通过 & $\checkmark$ \\
F12 & Merkle快照 & 管理员触发 & Root入审计日志 & 通过 & $\checkmark$ \\
F13 & 条件溯源 & 用户申请+审批 & 返回用户信息 & 通过 & $\checkmark$ \\
F14 & 三端一致性 & 管理员验证 & 三方Root一致 & 通过 & $\checkmark$ \\
\hline
\end{tabular}
\end{table}

其中，Token 生成采用$\mathrm{HMAC\text{-}SM3}(K_{\mathrm{order}}, \mathit{orderID} \| \mathrm{PID} \| \mathit{timestamp} \| \mathit{nonce})$，验证时后端从 session 获取 PID，不信任前端提交值。消息写入前经 SM4-CBC 加密（独立 IV），数据库存储$\mathrm{IV}+\mathrm{ciphertext}$。三端一致性验证比对管理员、用户和骑手各自上报的 Merkle Root，三者完全一致方判定通过。
